\begin{frame}[plain, noframenumbering]
	\begin{center}
		\Huge \textcolor{blue}{\textbf{Methodology}} \\
		\vspace{0.5cm}
		\large \textit{Data Processing, MD detections and Tracking}
	\end{center}
\end{frame}

\subsection{Dataset Information}
\begin{frame}{Dataset Information}
	\vspace{-0.1cm}
	\begin{block}{ERA5 Single Level Reanalysis Dataset (June-September 2024) \tiny \cite{hersbach2020era5} :}
		%\pdflink{2020-Hersbach-The-ERA5-global-reanalysis.pdf}{The ERA5 global reanalysis}
		\footnotesize
		% FIX: the spatial (31 km globally) and temporal (hourly) resolutions of ERA5
		\begin{itemize}
			\item \textbf{Spatial resolution :} \texttt{$0.25^\circ \times 0.25^\circ$}
			\item \textbf{Temporal resolution :} Hourly
			\item \textbf{Spatial Domain :} $10$--$30^\circ$ N, $70$--$95^\circ$ E (Indian Subcontinent)
			\item \textbf{Temporal domain :} June to September 2024 (4 months, Indian monsoon season)
		\end{itemize}
	\end{block}

	\vspace{-0.2cm}
	\begin{columns}[T, totalwidth=0.9\textwidth]
		\column{0.43\textwidth}
		\begin{block}{ERA5 Dataset Variables :}
			\begin{itemize}\footnotesize
				\item Mean Sea Level Pressure (\texttt{MSL})
				\item Surface Pressure (\texttt{SP})
				\item 10-Meter U-Wind Component (\texttt{U10})
				\item 10-Meter V-Wind Component (\texttt{V10})
				\item Mean Surface Latent Heat Flux (\texttt{AVG\_SLHF})
				\item Mean surface Sensible Heat Flux (\texttt{AVG\_ISHF})
				\item Total Precipitation (\texttt{TP})
				\item Volumetric Soil Water Layer-1 (\texttt{SWVL1})
				\item Soil Temperature Level-1 (\texttt{STL1})
				\item 2-Meter Temperature (\texttt{T2M})
				\item 2-Meter Dew Point Temperature (\texttt{D2M})
			\end{itemize}
		\end{block}

		\column{0.55\textwidth}
		\begin{block}{Derived Variables :}
			\begin{itemize}\footnotesize
				\item \textbf{10-meter Wind Speed (\texttt{W10}) :}
				      \footnotesize
				      \begin{equation}
					      \texttt{W10} = \sqrt{(\texttt{U10})^2 + (\texttt{V10})^2}
				      \end{equation}
				\item \textbf{Relative Vorticity (\texttt{$\zeta$}):} Calculated using the 10-m horizontal wind components
				      \footnotesize
				      \begin{equation}
					      \zeta = \frac{\partial(\texttt{V10)}}{\partial x} - \frac{\partial(\texttt{U10)}}{\partial y}
				      \end{equation}
				\item \textbf{Total Heat Flux (\texttt{THF}) :} Sum of surface latent and sensible heat fluxes
				      \footnotesize
				      \begin{equation}
					      \texttt{THF} = \texttt{SLHF} + \texttt{SSHF}
				      \end{equation}
				\item \textbf{Saturation Vapor Pressure (\texttt{$e_s$}) :}
				      \footnotesize
				      \begin{equation}
					      e_s = 6.112 \exp\left(\frac{17.67 * D2M}{D2M + 243.5}\right)
				      \end{equation}
				      where \texttt{D2M} = 2m Dew Point Temperature in $^\circ$C
				\item \textbf{Water Vapor Mixing Ratio (\texttt{w}) :}
				      \footnotesize
				      \begin{equation}
					      \texttt{w} = 0.622 \frac{e_s}{SP - e_s} \times 1000
				      \end{equation}
			\end{itemize}
		\end{block}
	\end{columns}
	% \item \textbf{Processing :} Multi-month concatenation, temporal regularization, spatial clipping
\end{frame}

\subsection{MD Detection Algorithm}
\begin{frame}{Monsoon Depression Detection Algorithm}
	\vspace{-0.1cm}
	\begin{columns}[T, onlytextwidth]
		\column{0.48\textwidth}
		\footnotesize
		\begin{block}{MD Detections Criteria :}
			\begin{itemize}
				\item \textbf{Pressure (\texttt{MSL}) :} \texttt{< 1002 hPa}
				\item \textbf{Contour Closure :} Closed isobar within \texttt{5$^{\circ}$} buffer
				\item \textbf{Track Duration :} \texttt{$\geq$ 24 hours}
				\item \textbf{Min net displacement:} \texttt{250 km}
				\item \textbf{Wind speed range\footnotemark[1] :} \texttt{5} to \texttt{19 m/s}
				\item \textbf{Total Precipitation (\texttt{TP}) :} \texttt{> 0.3 mm/hr} within buffer
				\item \textbf{Vorticity (\texttt{$\zeta$}) :} \texttt{> 0 /s} (cyclonic)
				\item \textbf{Max MD centre speed :} \texttt{12 m/s}
			\end{itemize}
		\end{block}

		\column{0.52\textwidth}
		\footnotesize
		\begin{block}{Tracking \& Post-Processing :}
			\begin{itemize}
				\item \textbf{Hourly Linking :} \texttt{50 km} search radius
				\item \textbf{Gap Bridging :} Up to \texttt{24 h}
				\item \textbf{Merge Overlaps :} \texttt{250 km} radius
				\item \textbf{Filtering :} WNW preference (\texttt{240$^{\circ}$ – 30$^{\circ}$})
				\item \textbf{Landfall Detection :}
				      \begin{itemize}
					      \footnotesize
					      \item \textbf{Database:} Natural Earth Physical Mask
					      \item \textbf{Resolution:} \texttt{50m} (1:50M scale)
					      \item \textbf{Method:} Point-in-polygon check
				      \end{itemize}
			\end{itemize}
		\end{block}

		% TODO: check SWVVL@ threshold
		\begin{block}{Analysis Parameters :}
			\begin{itemize}
				\item \textbf{\texttt{SWVL1} Range\footnotemark[2] :} \texttt{0.02 - 0.55 m$^3$/m$^3$}
				\item \textbf{Max Lag :} \texttt{$\pm$ 48 hours}
				\item \textbf{Long-lived Threshold :} \texttt{> 96 hours} (\texttt{> 4 days})
				\item \textbf{Statistical Significance :} \texttt{p < 0.05}
			\end{itemize}
		\end{block}
	\end{columns}

	\vfill
	\footnotetext[1]{\tiny Wind speed criteria based on IMD Monsoon Report 2024 (page 284), Source : \href{https://mausam.imd.gov.in/imd\_latest/monsoonreport2024.pdf}{\url{https://mausam.imd.gov.in/imd\_latest/monsoonreport2024.pdf}}}
	\footnotetext[2]{\tiny SWVL1 range based on ERA5 climatology for Indian region, Source : \href{https://www.researchgate.net/publication/365632311\_Assessment\_of\_ERA5-Land\_volumetric\_soil\_water\_layer\_product\_using\_in-situ\_and\_SMAP\_soil\_moisture\_observations}{\url{https://www.researchgate.net/publication/365632311\_Assessment\_of\_ERA5-Land\_volumetric\_soil\_water\_layer\_product\_using\_in-situ\_and\_SMAP\_soil\_moisture\_observations}}}
\end{frame}

% % WARN: check this before adding otherwise remove it
% % \begin{frame}{Track Detection Algorithm}
% %     \begin{block}{Vortex Detection (Okubo-Weiss Parameter)}
% %         \begin{enumerate}
% %             \item Extract relative vorticity and deformation rate from ERA5 wind fields
% %             \item Compute Okubo-Weiss parameter : $OW = \delta^2 + D^2 - \zeta^2$
% %             \item Identify closed circulation centers where $OW < -10^{-10}$ s$^{-2}$
% %             \item Track temporal continuity using Munkres assignment algorithm
% %         \end{enumerate}
% %     \end{block}
% % \end{frame}
% % \begin{frame}{Analysis Pipeline}
% %     \begin{block}{Processing Stages}
% %         \begin{enumerate}
% %             \item \textbf{Data Loading :} Combine monthly netCDF files, regularize to \verb|1-hour| resolution
% %             \item \textbf{Detection :} Apply MD criteria hourly, identify low-pressure centers
% %             \item \textbf{Tracking :} Link detections, stitch gaps, merge overlaps, smooth tracks
% %             \item \textbf{Per-Track Analysis :}
% %             \begin{itemize}
% %                 \item Translational speed calculation
% %                 \item Azimuthal Hovm\"oller diagrams (\verb|25–400 km| radii)
% %                 \item Soil moisture–precipitation correlation (\verb|24-hour| lags)
% %                 \item Surface heat flux evolution (\verb|latent/sensible|)
% %                 \item Landfall transition analysis (\verb|24 hours| before/after)
% %             \end{itemize}
% %             \item \textbf{Global Analysis :}
% %             \begin{itemize}
% %                 \item Coupling metrics (\verb|FSI|, diurnal amplitude, efficiency)
% %                 \item Antecedent soil moisture patterns
% %                 \item Classification (\verb|short-lived| vs. \verb|long-lived| MDs)
% %             \end{itemize}
% %             \item \textbf{output\_2024 :} Maps, plots, CSV exports, summary report
% %         \end{enumerate}
% %     \end{block}
% % \end{frame}
%
% % TODO: Verify from correlation_analysis.py
% \subsection{Correlation Analysis}
% \begin{frame}{Correlation Analysis}
%     Quantifies coupling strength : Does increase in rainfall lead to soil moisture increment?
%
%     \vspace{0.2cm}
%     Positive coupling $\Rightarrow$ Self-sustaining feedback possible \\
%     Negative or weak coupling $\Rightarrow$ No reinforcing feedback
%
%     \vspace{0.2cm}
%     \footnotesize
%     For each individual track point :
%     \begin{itemize}
%         \footnotesize
%         \item \textbf{Data Extraction :} Extract \texttt{TP} (Total Precipitation, Accumulated hourly) and \texttt{SWVL1} (Soil Water Volumetric Level-1, \texttt{0-7 cm}).
%         \item \textbf{Spatial Buffering :} Apply a \texttt{100 km} radius buffer around the MD center to capture regional interaction.
%         \item \textbf{Regional Averaging :} Compute the mean value within the buffer at each time step to reduce grid-scale noise.
%         \item \textbf{Statistical Processing :}
%         \begin{itemize}
%             \footnotesize
%             \item Calculate Pearson correlation (\texttt{r}) across lags : \texttt{0h, 6h, 12h, 24h, 48h}.
%             \item Adjust significance using \textbf{Effective Sample Size} to account for temporal autocorrelation.
%             \item Output : Identify the \textbf{Peak Lag} where \texttt{|r|} is maximized.
%         \end{itemize}
%         \item Pearson correlation with t-test significance.
%         \item Linear regression for trend analysis.
%         \item Lagged analysis for temporal relationships.
%         \item Spatial interpolation (nearest-neighbor).
%         % TODO:
%         % \item Quality Assurance:
%         % \begin{itemize}
%         %     \footnotesize
%         %     \item Outlier detection and documentation
%         %     \item Cross-validation with satellite observations
%         %     \item Sanity checks on derived quantities
%         %     \item Uncertainty quantification where possible
%         % \end{itemize}
%     \end{itemize}
% \end{frame}
%
% \subsection{Lag Correlation}
% % \begin{frame}{Methodology: Lag Correlation \& Interpretation}
% \begin{frame}{Lag Correlation}
%     \vspace{0.2cm}
%     Lag correlation identifies the temporal sequencing of physical processes, determining which variable acts as the "driver" and which as the "responder."
%
%     \vspace{0.15cm}
%     \footnotesize
%     \begin{itemize}
%         \item \textbf{Lag (Time Offset) :} The time shift (in \texttt{hours}) applied to one variable relative to another (\texttt{TP} relative to \texttt{SWVL1}).
%         \item \textbf{Leading Variable (\textit{Driver}) :} The variable shifting first. If \texttt{TP} leads \texttt{SWVL1}, rain precedes soil moisture increase.
%         \item \textbf{Lagging Variable (\textit{Responder}) :} The variable responding to the initial driver forcing.
%         \item \textbf{Peak Correlation (\texttt{r$_{max}$}) :} The lag hour where the relationship (\texttt{|r|}) is strongest, representing the system's "reaction time".
%     \end{itemize}
%     % \begin{itemize}
%     %     \item \textbf{Lag (Time Offset):} The time shift in \texttt{hours} applied to \texttt{TP} relative to \texttt{SWVL1}.
%     %     \item \textbf{Positive Lag ($+h$):} \texttt{TP} leads \texttt{SWVL1}. This confirms a physical "Handshake" where rainfall recharges the soil.
%     %     \item \textbf{Peak Correlation ($r_{max}$):} The lag hour where the relationship (\texttt{|r|}) is strongest/maximum, representing the system's \textbf{infiltration reaction time}.
%     % \end{itemize}
%
%     \vspace{0.2cm}
%     \begin{block}{Interpretation of the Plots}
%         \begin{itemize}
%             \item \textbf{X-axis (\texttt{Lag}) :} Time shift (\texttt{hours}) between the two variables.
%             \item \textbf{Y-axis (Correlation \texttt{r}) :} Relationship strength (\texttt{-1.0} to \texttt{1.0}).
%             \item \textbf{Significance :} Indicated by \texttt{p-values < 0.05} based on effective degrees of freedom.
%             \item \textbf{The Infiltration Peak :} For lag-correlation plot for \texttt{TP-SWVL1} if clear peak at \texttt{+3} to \texttt{+6h} indicates efficient moisture capture. A flat or zero-lag line suggests \textbf{Rapid Runoff} without storage.
%             \item \textbf{The Saturation Plateau :} If correlations collapse while SM remains high ($>0.40$), the soil has reached capacity. Additional \texttt{TP} no longer drives \texttt{SWVL1} but triggers \textbf{Surface Runoff}.
%             \item \textbf{The Memory (\texttt{ACF}) Decay :} A broad, slowly decaying correlation curve indicates high \textbf{Soil Memory} (Persistence), vital for the "Brown Ocean" effect. \cite{seneviratne2010investigating}
%         \end{itemize}
%     \end{block}
%
%     \vspace{0.15cm}
%     \footnotesize
%     \textbf{Significance:} Calculations use \texttt{Pearson r} with \texttt{p-values < 0.05} to ensure the feedback is not due to random noise.
% \end{frame}
%
% \subsection{Coupling Metrics}
% \begin{frame}{Coupling Metrics: Quantifying Land-Atmosphere Interactions}
%     \begin{block}{5 Key Metrics for Analysis:}
%         \begin{enumerate}
%             \item Feedback Strength Index (\texttt{FSI}, \texttt{0-1}) % Correlation strength + Capacity index
%             \item Coupling Efficiency (\texttt{CE})
%             \item Diurnal Sensible Heat Flux Amplitude ($A_{\texttt{SHF}}$)
%             \item Soil Moisture Memory Index (\texttt{ACF} of \texttt{SWVL1} at Lag \texttt{24hr})
%             \item Topographic Sensitivity Index (\texttt{TSI}, \texttt{0-1}) % MD behavior + terrain characteristics
%             % TODO: future work
%             % \item NDVI = Normalized Difference Vegetation Index = (NIR - RED) / (NIR + RED) -> Derived from: Red and near-infrared reflectances
%             % \item CAPE = Convective Available Potential Energy (J/kg) -> Derived from: Temperature, humidity, pressure profiles
%             % \item Drought Index = f(rainfall, temperature, soil moisture, ET) -> Derived from: Multiple raw parameters
%         \end{enumerate}
%     \end{block}
%
%     \vspace{0.4cm}
%     \begin{center}
%     \small \textit{All five metrics work together to define the system's state:}
%     % Feedback Strength Index (FSI)
%     %            ↓ (component of)
%     %   Coupling Efficiency (CE)    <- Diurnal Heat Flux Amplitude (A_SHF)
%     %           ↓                                ↓
%     % Topographic Sensitivity (TSI) <- Capacity Index
%     %           ↓
%     % Soil Moisture Memory (ACF)
%     %     (parallel, independent measurement)
%     \[
%         \underbrace{\mathtt{FSI}}_{\substack{\text{Correlation,} \\ \text{Coupling Strength}}}
%         + \underbrace{\mathtt{A_{SHF}}}_{\substack{\text{Surface Heat} \\ \text{Flux Amplitude} \\ \text{Tells Response}}}
%         + \underbrace{\mathtt{ACF_{24h}}}_{\substack{\text{Soil Memory} \\ \text{Water holding capacity} \\ \text{Determines Persistence}}}
%         + \underbrace{\mathtt{TSI}}_{\substack{\text{Terrain} \\ \text{Effect}}}
%         \Rightarrow \underbrace{\mathtt{CE}}_{\substack{\text{Coupling} \\ \text{Efficiency}}}
%     \]
%     \end{center}
% \end{frame}
%
% \subsubsection{Feedback Strength Index (\texttt{FSI})}
% % Formula:    FSI = |Corr(TP, SWVL1)| x Capacity Index x Significance Factor
% % Range:      [0, 1]
% % Parameters: 
% %   - SWVL1_dry   = 0.02 m³/m³ (wilting point)
% %   - SWVL1_wet   = 0.55 m³/m³ (saturation point)
% %   - p_threshold = 0.05
% %   - sig_factor  = 1.0 if p<0.05, else 0.5
% % Code:       coupling_metrics.py, lines 41-76, [compute_feedback_strength_index() lines 41-76], [compute_soil_water_capacity_index() lines 28-39]
% % FSI data:           coupling_metrics_FSI.csv
% % 2024 Mean:  0.502 (correlation strength), 0.530 (capacity)
% % Status:     4 of 6 tracks significant (p<0.05)
% % FSI 0.0 - 0.2: Weak coupling (insufficient feedback)
% % FSI 0.2 - 0.4: Moderate coupling (some feedback)
% % FSI 0.4 - 1.0: Strong coupling (robust feedback)
% \begin{frame}[label=feedbackstrengthindex]{Coupling Metrics [1/5]: Feedback Strength Index (\texttt{FSI})}
%     \vspace{-0.3cm}
%     \footnotesize
%     \begin{block}{1. Feedback Strength Index (\texttt{FSI})}
%         \textbf{Mathematical Definition:}
%         \begin{equation}
%             \text{\texttt{FSI}} = |\text{Corr}(\text{\texttt{TP, SWVL1}})| \times \text{\texttt{Capacity Index}} \times \text{\texttt{Significance Factor}}
%             \label{eq:fsi}
%         \end{equation}
%
%         \textbf{Parameters :}
%         \begin{itemize}
%             \item \textbf{\texttt{Capacity Index} (Soil water capacity index)} = $\dfrac{\text{mean}(\text{\texttt{SWVL1}}) - \texttt{0.02}}{\texttt{0.55} - \texttt{0.02}}$, clipped to $[\texttt{0,1}]$
%             \begin{itemize}
%                 \footnotesize
%                 \item Available soil moisture capacity relative to saturation level
%                 \item \texttt{0.02 m$^3$/m$^3$} = Wilting point (dry soil, ERA5 climatology)\footnotemark[3]
%                 \item \texttt{0.55 m$^3$/m$^3$} = Saturation point (wet soil, ERA5 climatology)\footnotemark[3]
%             \end{itemize}
%             \item \textbf{\texttt{Significance Factor}} = \texttt{1.0} if \texttt{p < 0.05}, else \texttt{0.5}
%         \end{itemize}
%
%         Integrates three factors determining coupling sustainability :
%         \begin{enumerate}
%             \item \textbf{Correlation Strength :} How tightly \texttt{TP} links to \texttt{SWVL1}. Range: \texttt{0} (no feedback) to \texttt{1} (perfect feedback).
%             \item \textbf{Soil Water Capacity Index :} The soil's ability to hold/release moisture based on type (\texttt{0} = dry, \texttt{1} = saturated).
%             \item \textbf{Statistical Significance :} Robustness of the link (\texttt{p < 0.05}). Factor of \texttt{1.0} if significant, \texttt{0.5} otherwise.
%         \end{enumerate}
%     \end{block}
%
%     \footnotetext[3]{\tiny \texttt{SWVL1} range based on ERA5 climatology for Indian region, Source : \href{https://www.researchgate.net/publication/365632311\_Assessment\_of\_ERA5-Land\_volumetric\_soil\_water\_layer\_product\_using\_in-situ\_and\_SMAP\_soil\_moisture\_observations}{\url{https://www.researchgate.net/publication/365632311\_Assessment\_of\_ERA5-Land\_volumetric\_soil\_water\_layer\_product\_using\_in-situ\_and\_SMAP\_soil\_moisture\_observations}}}
% \end{frame}
%
% \subsubsection{Coupling Efficiency (\texttt{CE})}
% % Formula:    CE = (FSI x A_norm) / Duration_Factor
% % Range:      [0, ∞) typically [0.5, 15]
% % Parameters:
% %   - A_norm           = tanh(A_SHF / 100)  [saturation curve]
% %   - Duration_Factor  = max(1.0, Hours / 48)
% %   - Reference point  = 48 hours
% % Code:       coupling_metrics.py, lines 264-292 [compute_coupling_efficiency() (lines 264-292)]
% % CE data:            coupling_metrics_efficiency.csv
% % 2024 Mean:  8.76 \pm 8.80
% % 2024 Range: 0.5 - 15+
% % All Tracks: CE > 0.5 (above minimum threshold)
% % CE : 0.5 - 3.0: Weak coupling (limited persistence)
% % CE : 3.0 - 8.0: Moderate coupling (standard persistence)
% % CE : 8.0 - 15+: Strong coupling (self-sustaining)
% \begin{frame}[label=couplingefficiency]{Coupling Metrics [2/5]: Coupling Efficiency (\texttt{CE})}
%     \vspace{-0.1cm}
%     \footnotesize
%     \begin{block}{2. Coupling Efficiency (\texttt{CE})}
%         Quantifies the energy exchange intensity relative to the system's lifespan.
%
%         \textbf{Mathematical Definition :}
%         \begin{equation}
%             \text{\texttt{CE}} = \frac{\text{\texttt{FSI}} \times A_{\text{\texttt{norm}}}}{\text{\texttt{Duration Factor}}}
%             \label{eq:couplingeff}
%         \end{equation}
%
%         \textbf{Parameters :}
%         \begin{itemize}
%             \item \textbf{\texttt{FSI}} = Feedback Strength Index $\in [\texttt{0, 1}]$ (from equation \ref{eq:fsi})
%             \item $A_{\text{\texttt{norm}}} = \tanh(A_{\text{\texttt{SHF}}} / \texttt{100})$ = Normalized amplitude via saturation \texttt{tanh} function (range: $[\texttt{0, 1}]$)
%             \item \textbf{\texttt{Duration Factor}} = $\max(\texttt{1.0}, \text{\texttt{Duration}}_{\text{\texttt{hours}}} / \texttt{48})$ (normalizes systems $>\texttt{48 hours}$)
%         \end{itemize}
%
%         \begin{enumerate}
%             \item \texttt{FSI} : Measures strength of rainfall-to-soil moisture (\texttt{TP-SWVL1}) feedback mechanism.
%             \item \texttt{Normalized Amplitude for SHF (\texttt{A$_{SHF}$})} : Captures intensity of surface response (saturating function prevents outliers).
%             \item \texttt{Duration Factor} : Prevents long-lived systems from dominating purely by time accumulation; systems $>\texttt{48h}$ show proportionally weaker coupling signal.
%         \end{enumerate}
%
%         \vspace{0.2cm}
%         \textbf{Interpretation :}
%         \begin{itemize}
%             \item \textbf{High \texttt{CE}} ($\texttt{8-15}$) = Efficient energy exchange $\Rightarrow$ Self-sustaining.
%             \item \textbf{Moderate \texttt{CE}} ($\texttt{3-8}$) = Balanced coupling with some duration effects.
%             \item \textbf{Low \texttt{CE}} ($<\texttt{3}$) = Weak coupling $\Rightarrow$ Limited inland persistence.
%         \end{itemize}
%
%         \vspace{0.3cm}
%         \footnotesize
%         \textbf{Note :} \texttt{FSI} Discriminates between: (A) Energy-driven survival (\textbf{High \texttt{CE}}) vs. (B) Moisture-managed survival fur to high soil moisture and precipitation (\textbf{Low \texttt{CE}} + \textbf{High \texttt{ACF$_{24h}$}})
%     \end{block}
% \end{frame}
%
% \subsubsection{Diurnal \texttt{SHF} Amplitude}% ($A_{\text{\texttt{SHF}}}$)  }
% % Formula:    A_SHF = SHF_max - SHF_min  (land phase only)
% % Units:      W/m$^2$
% % Code:       coupling_metrics.py, lines 120-227, [compute_diurnal_shf_amplitude() (lines 121-171): Computes amplitude for given phase], [compute_diurnal_amplitude_by_phase() (lines 173-227): Separates ocean vs land], [compute_diurnal_amplitude_all_tracks() (lines 228-262): Applied to all tracks]
% % A_SHF data:         coupling_metrics_diurnal_amplitude.csv
% % 2024 Values (Mean \pm Range):
% %   Ocean Phase:  13.54 W/m$^2$ (mean: -11.44 W/m$^2$)
% %   Land Phase:   77.57 W/m$^2$ (mean: +2.08 W/m$^2$)
% %   Amplitude Shift: ΔA = 64.03 W/m$^2$ (~5.7x)
% % Physical:   SHF = ρCp u*(T_s - T_a)
% %            Land heats -> positive SHF -> CAPE enhancement
% %            Ocean stable -> negative SHF -> suppressed convection
% % A_SHF : < 30 W/m$^2$: Minimal land impact
% % A_SHF : 30 - 50 W/m$^2$: Moderate land heating
% % A_SHF : > 50 W/m$^2$: Strong land heating
% \begin{frame}[label=diurnal]{Coupling Metrics [3/5]: Diurnal \texttt{SHF} Amplitude ($A_{\text{\texttt{SHF}}}$)}
%     \footnotesize
%     \begin{block}{3. Diurnal Sensible Heat Flux Amplitude ($A_{\text{\texttt{SHF}}}$)}
%         \textbf{Mathematical Definition :}
%         \begin{equation}
%             A_{\text{\texttt{SHF}}} = \text{\texttt{SHF}}_{\text{\texttt{max}}} - \text{\texttt{SHF}}_{\text{\texttt{min}}} \quad \text{(Calculated over Land Phase)}
%         \end{equation}
%
%         \begin{itemize}
%             \item \textbf{Ocean phase :} $\lesssim \texttt{15}$ (Damped by thermal inertia).
%             \item \textbf{Land phase :} $\gtrsim \texttt{65}$ (Strong daytime heating).
%             \item \textbf{Amplitude Shift :} $\Delta A \approx \texttt{40}$ for long-lived MDs.
%         \end{itemize}
%     \end{block}
%
%     \vspace{0.2cm}
%     \textbf{Physical Mechanism :}
%     \vspace{-0.1cm}
%     \begin{equation}
%         \text{Sensible Heat Flux \texttt{SHF}} = \rho C_p u_* (T_s - T_a)
%         \label{eq:shf}
%     \end{equation}
%
%     \begin{itemize}
%         \item Measures land surface's diurnal heating response. Over land, surface temperature varies with solar radiation (strong diurnal cycle).
%         \item \textbf{Ocean :} Water's thermal inertia suppresses the cycle $\rightarrow$ \texttt{SHF} nearly constant.
%         \item \textbf{Land :} Soil heats/cools rapidly $\rightarrow$ Large \texttt{SHF} swings.
%         \item Large amplitude $\Rightarrow$ Stronger destabilization fueling deep convection.
%     \end{itemize}
%
%     \vspace{0.2cm}
%     \textbf{Interpretation :} Measures land surface's diurnal heating response
%
%     \begin{block}{Computation :}
%         \begin{enumerate}
%             \item Extracted hourly \texttt{SHF} from \texttt{ERA5} along track.
%             \item Classify points using \texttt{is\_land} mask (\texttt{True/False}).
%             \item Calculate $A_{\text{\texttt{SHF}}}$  $=(Max(\text{\texttt{SHF}}) - Min(\text{\texttt{SHF}}))$ for land and ocean phase.
%             \item Determine $\Delta A = A_{\text{\texttt{SHF,land}}} - A_{\text{\texttt{SHF,ocean}}}$.
%         \end{enumerate}
%     \end{block}
% \end{frame}
%
% \subsubsection{Soil Moisture Memory (\texttt{ACF})}
% % Formula:    R(τ) = Cov(X_t, X_t+τ) / σ$^2$
% % Lag:        24 hours (filters diurnal cycle)
% % Range:      [-1, 1]
% % Code:       soilmoiture_analysis.py, lines 23-42
% % ACF data:           soil_memory_analysis.csv
% % Library:    statsmodels.tsa.stattools.acf (FFT-based)
% % Thresholds:
% %   ACF : > 0.35: High memory (multi-day persistence)   (clay-rich, retains water)
% %   ACF : 0.1 - 0.35: Moderate memory (some persistence)
% %   ACF : < 0.1: Low memory (isolated events)   (sandy, drains quickly)
% % Mechanism:  Measures soil's ability to sustain feedback
% %            High ACF -> Longer feedback persistence
% %            Low ACF  -> Isolated rainfall events
% \begin{frame}[label=soilmoiturememory]{Coupling Metrics [4/5]: Soil Moisture Memory (\texttt{ACF})}
%     \vspace{-0.2cm}
%     \footnotesize
%     \begin{block}{4. Soil Moisture Memory Index (\texttt{ACF} at \texttt{24-hour} Lag)}
%         Persistence of soil moisture anomalies after the initial \texttt{TP} forcing has ceased.
%
%         \vspace{0.2cm}
%         \textbf{Mathematical Definition : Autocorrelation function (\texttt{ACF}):} \\
%         Correlation of variable with itself 24 hours prior.
%         \begin{equation}
%             R(\tau) = \frac{E[(X_t - \mu)(X_{t+\tau} - \mu)]}{\sigma^2}
%         \end{equation}
%         High \texttt{ACF} = High carry over of latent heat potential.
%
%         \vspace{0.2cm}
%         \textbf{Memory Index at \texttt{24h} lag (\texttt{ACF$_{24h}$}) :}
%         \begin{equation}
%             \text{\texttt{ACF}}(\text{\texttt{SWVL1}}, \text{lag}=\texttt{24h}) = \frac{\text{Cov}(\text{\texttt{SWVL1}}_t, \text{\texttt{SWVL1}}_{t+24})}{\sigma_{\text{\texttt{SWVL1}}}^2}
%         \end{equation} \cite{seneviratne2010investigating}
%
%         \textbf{Lag Choice :} \texttt{24-hour} lag filters out the rapid diurnal cycle (daytime evaporation+transpiration(ET) Vs. nighttime recovery) to measure real multi-day moisture persistence.% Computed via \texttt{statsmodels.acf} with FFT.
%     \end{block}
%
%     \vspace{0.2cm}
%     \footnotesize
%     \textbf{Interpretation :}
%         \begin{itemize}
%             \item \textbf{High Memory ($> \texttt{0.35}$):} \texttt{SWVL1} today strongly predicts tomorrow's moisture content.
%                 \footnotesize
%                 \begin{itemize}
%                     \item Soil retains water effectively (e.g., clay-rich, high \texttt{Capacity Index}).(\tiny Only Clay has the capillary strength to multi-day retention \footnotesize)
%                     \item Enables sustained feedback : rainfall $\rightarrow$ soil saturation $\rightarrow$ enhanced ET $\rightarrow$ enhanced convection.
%                 \end{itemize}
%             \item \textbf{Low/Negative Memory ($< \texttt{0.1}$):} \texttt{SWVL1} is uncorrelated with the next day. Soil rapidly loses moisture.
%                 \footnotesize
%                 \begin{itemize}
%                     \item Soil and rocky terrain drains quickly (e.g., sandy, low \texttt{Capacity Index}). Sand lacks capillary retention.
%                     \item Moisture removed by rapid evaporation, drainage, plant uptake or other natural processes.
%                     \item Weak feedback potential; each rainfall is isolated, not sustained.
%                 \end{itemize}
%         \end{itemize}
% \end{frame}
%
% \subsubsection{Terrain Strength Index (\texttt{TSI})}
% % Formula:    TSI = FSI x 0.6 + Capacity Index x 0.4
% % Range:      [0, 1]
% % Weighting:
% %   - FSI (0.6):      Precipitation-feedback mechanism [dominates]
% %   - Capacity (0.4): Soil type distribution effects [modulates]
% % Code:       statistics_analysis.py, lines 291-360, [compute_tsi_index() (lines 291-360)]
% % TSI data:           h7_tsi_distribution.png (histogram)
% % 2024 Mean:  0.39 \pm 0.21
% % 2024 Range: 0.179 - 0.713
% % Interpretation: Terrain explains ~39% of depression behavior
% %                Other factors (SST, wind, humidity): ~61%
% % Physics:    - Orographic lifting -> enhanced windward precip
% %             - Barrier effects -> flow modification  
% %             - Differential heating -> boundary layer changes
% %             - Soil heterogeneity -> variable moisture
% % TSI : 0.0 - 0.2: Terrain negligible
% % TSI : 0.2 - 0.4: Terrain moderate (highest for 2024: 0.39)
% % TSI : 0.4 - 0.6: Terrain significant
% % TSI : 0.6 - 1.0: Terrain dominant
% % 0.6 and 0.4 weighting in TSI -> FSI (0.6) dominates because it captures the actual precipitation-to-soil moisture feedback that terrain influences. Capacity Index (0.4) modulates because terrain affects soil type distribution. Code: in statistics_analysis.py
% % \begin{frame}[label=terrainstrenghtindex]{Coupling Metrics [5/5]: Terrain Strength Index (\texttt{TSI})}
% \begin{frame}[label=terrainstrenghtindex]{Coupling Metrics [5/5]: Terrain Strength Index (\texttt{TSI})}
%     \vspace{-0.2cm}
%     \footnotesize
%     \begin{block}{5. Terrain Strength Index (\texttt{TSI})}
% %       % a dimensionless derived metric that quantifies the proportion of monsoon depression behavior variance attributable to topographic and elevation effects.
%       A dimensionless metric that quantifies how much of the depression's behavior variance is influenced by terrain characteristics (as opposed to ocean-driven or atmospheric factors).
%
%       \vspace{0.2cm}
%       % TODO:
%       % TSI = 1 - R$^2$(MD_metrics | without_terrain_predictors) = [Var_explained_by_terrain] / [Total_variance_in_MD]
%       \textbf{Mathematical Definition :}
%       % \text{\texttt{TSI}} = \texttt{FSI} \times 0.6 + \text{\texttt{Capacity Index}} \times 0.4
%       \begin{equation}
%            \mathrm{TSI} = \mathrm{FSI} \times 0.6 + \mathrm{Capacity~Index} \times 0.4
%       \end{equation}
%       {\tiny where \texttt{TSI} makes link between feedback strength (60\%) with soil water capacity (40\%) to estimate terrain-modulated response}
%       % {\tiny where MD_metrics include intensity evolution, precipitation rate, track velocity, and lifetime duration.} TODO:  \cite{}
%
%       \textbf{Rationale :}
%       \begin{itemize}
%          \footnotesize
%          \item \textbf{Weighting:} Feedback strength dominates (\texttt{60\%}) because terrain affects depression dynamics through soil moisture feedback (orographic lift $\rightarrow$ rainfall $\rightarrow$ wet soil).
%          \item \textbf{Capacity component:} Terrain constrains soil type (e.g., clay in valleys, sand on ridges), influencing water retention.
%          \item \textbf{Range:} \texttt{0} $\rightarrow$ \texttt{1}, where higher values indicate terrain's strong influence on system behavior.
%       \end{itemize}
%    \end{block}
%
%    \vspace{0.15cm}
%    \textbf{Physical Mechanisms:}
%    \footnotesize
%    \begin{itemize}
%        \item \textbf{Orographic lifting:} Mountains block/redirect airflow, enhance precipitation on windward slopes.
%        \item \textbf{Barrier/channeling:} Terrain confines and accelerates flow, altering wind structure and pressure fields.
%        \item \textbf{Differential heating:} Exposed slopes heat differently than valleys; affects boundary layer instability.
%        \item \textbf{Soil heterogeneity:} Mountainous regions have patchy soil types (clay pockets, sandy ridges) $\rightarrow$ variable moisture response.
%    \end{itemize}
%
%    \vspace{0.15cm}
%    \textbf{Interpretation :}
%    \footnotesize
%      \begin{itemize}
%            \item TSI = 0 : No terrain sensitivity, depression is fully driven by other factors like wind, humidity, large-scale monsoon circulation etc.
%            \item 0 < TSI < 1 : Partial terrain sensitivity
%            \item TSI = 1 : Complete terrain determination
%            \item \textbf{Mean TSI :} Fraction of depression behavior due to terrain for all tracks.
%            \item \textbf{Std. TSI :} Variation of how different depressions have different  terrain sensitivities depending on track orientation relative to mountains, depression intensity, speed.
%      \end{itemize}
% %   TODO:
% %   Computation : 
% %      1. Extract elevation along MD track (SRTM 30-m)
% %      2. Compute terrain metrics: slope, aspect, relief, curvature
% %      3. Build statistical model: MD_behavior ~ f(terrain_predictors)
% %      4. Calculate R$^2$ (coefficient of determination)
% %      5. TSI = 1 - R$^2$
% %    This provides a quantitative measure of the importance of terrain effects for regional MD forecasting applications.
% \end{frame}
%
% % TODO: Add SSI (Standardized Soil-moisture Index)
%
% \subsection{Heat Flux Transition}
% \begin{frame}{Heat Flux Transition \& Brown Ocean Effect}
%     \footnotesize
%     \begin{block}{The Brown Ocean Effect (BoE) Hypothesis}
%         Typically applied to Tropical Cyclones; here applied to \textbf{Monsoon Depressions (MDs)} :
%         \begin{itemize}
%             \item \textbf{Core Concept :} Wet/saturated soil acts as a "pseudo-ocean," providing latent heat to sustain the system inland.
%             \item \textbf{MD Context :} Unlike TCs, MDs are already "moisture-rich" but weak. BoE suggests that high \texttt{SWVL1} maintains the depression's core structure by preventing the standard post-landfall energy collapse.
%         \end{itemize}
%     \end{block}
%
%     \begin{block}{Heat Flux Signatures at Landfall}
%         \begin{enumerate}
%             \item \textbf{Ocean Phase :} \texttt{LHF} dominates, \texttt{SHF} remains weak.
%             \footnotesize
%             \begin{itemize}
%                 \item Ocean evaporation is the primary thermodynamic fuel source.
%             \end{itemize}
%             \item \textbf{Land Phase :} \texttt{LHF} drops, but \texttt{SHF} increases due to surface heating.
%             \begin{itemize}
%                 \footnotesize
%                 \item \textbf{Positive \texttt{SHF} :} Diurnal land heating sustains atmospheric instability.
%                 \item \textbf{Negative \texttt{SHF} :} Surface act as an energy sink $\rightarrow$ Rapid decay.
%             \end{itemize}
%         \end{enumerate}
%     \end{block}
%
%     \begin{block}{Thresholds for MD Persistence}
%         \begin{itemize}
%             \item \textbf{Pre-landfall \texttt{LHF} :} \texttt{-120} to \texttt{-160 W/m$^2$} (High wind-driven evaporation, Infinite moisture supply by ocean).
%             \item \textbf{Post-landfall \texttt{LHF} :} \texttt{>-80 W/m$^2$} (and require \texttt{SWVL1 > 0.5} to mimic ocean).
%             \item \textbf{Land \texttt{SHF} stays same :} \texttt{$\approx$ +4.5 W/m$^2$} (Transition/shift from ocean thermal inertia to rapid diurnal soil heating).
%         \end{itemize}
%     \end{block}
% \end{frame}
%
% \subsection{\texttt{MODIS} Satellite Data}
% \begin{frame}{\texttt{MODIS} Data (\texttt{2000m} Scale) for ERA5 Tracks}
%     \begin{center}
%         \vspace{-0.2cm}
%         \footnotesize
%         \begin{block}{Data Acquisition \& GEE Implementation :}
%             \begin{itemize} \footnotesize
%                 \item \textbf{Platform :} Cloud-based geospatial analysis performed via \href{https://earthengine.google.com/}{Google Earth Engine (GEE)}\footnotemark[6] using the \textbf{JavaScript API}.
%                 \item \textbf{\texttt{MODIS} Dataset :} \texttt{MOD09GA.061} (Terra Surface Reflectance Daily \texttt{500m}).
%                 \item \textbf{Spectral Visualization :} True Color Composite using:
%                 \begin{itemize} \footnotesize
%                     \footnotesize
%                     \item \textbf{Red :} \texttt{sur\_refl\_b01}, \textbf{Green :} \texttt{sur\_refl\_b04}, \textbf{Blue :} \texttt{sur\_refl\_b03}
%                     \item \textbf{Gamma :} \texttt{1.3}, \textbf{Range :} \texttt{[0, 10000]}
%                 \end{itemize}
%                 \item \textbf{Temporal Processing :} Image collection filtered by \texttt{startDate} and \texttt{endDate} for each track; reduced using \texttt{.mean()} to create a cloud-persistent composite.
%                 \item \textbf{Spatial Integration :}
%                 \begin{itemize} \footnotesize
%                     \footnotesize
%                     \item \textbf{Vector Overlay :} ERA5-derived \texttt{track\_id} coordinates imported as an \texttt{ee.FeatureCollection} (Asset-hosted)\footnotemark[7].
%                     \item \textbf{Geometry :} Points converted to \texttt{ee.Geometry.LineString} with a \texttt{2000m} buffer for visibility.
%                     \item \textbf{Scale \& Projection :} Exported at \texttt{2000m} resolution (EPSG:\texttt{4326}).
%                 \end{itemize}
%             \end{itemize}
%         \end{block}
%
%         \vspace{0.3cm}
%         \begin{columns}[T]
%             \column{0.95\textwidth}
%                 \tiny *Note: The \texttt{MODIS} mean composite highlights the Central Dense Overcast (CDO). The color-coded tracks validate the spatial consistency between ERA5 low-pressure centers and real-world cloud mass.
%         \end{columns}
%     \end{center}
%
%     \vspace{-0.6cm}
%     \footnotetext[6]{\tiny Google Earth Engine (GEE) : \href{https://earthengine.google.com/}{\url{https://earthengine.google.com/}}}
%     \footnotetext[7]{\tiny GEE Implementation : \texttt{\href{https://code.earthengine.google.com/?asset=projects/final-year-project-466515/assets/MD-coordinates-ALL}{\url{projects/final-year-project-466515/assets/MD-coordinates-ALL}}}} %(for now it's private/closed-source)
% \end{frame}
